\section{Conclusions}\label{sec:conclusion}

This dissertation asked what a fiscal authority can learn from LLM-based exposure measures, given that they are demonstrably model-dependent. What it can learn is the differential, read scale-free instead of in the units of any one model's index. The levels themselves carry little. Scotland sits about two percentile points below the rest of the UK in the UK occupational exposure distribution, and 0.86 index points on the baseline model's 0--100 scale. The composed 95\% interval on that index-point figure, $[-1.07, -0.69]$, draws most of its width from survey sampling, and it describes precision within one model while saying nothing about agreement across models. London accounts for about two-fifths of that gap; against the UK excluding London, Scotland ranks fifth of the twelve regions and nations, slightly below the middle of a compressed distribution. The composition of Scottish exposure tilts towards substitution-dominated occupations even as both channel intensities sit below the rUK. Carried into the Acemoglu-style projection, the gap implies a ten-year TFP differential of $-0.029$ percentage points, and the stylised translation of \cref{sec:fiscal} converts that into a net-position sensitivity of about one and a half million pounds a year. That is modest against the \pounds969 million forecast net position, and well inside the SFC's own income tax forecast-error ranges, but it is persistent, and it is applied afresh at each three-year reconciliation instead of averaging out across them.

The contributions close as they opened. First, the robustness of the differential is established analytically and then verified, not asserted. Proposition 1 removes the common component of any model-specific bias and Corollary 1 bounds what survives, while Corollary 2 covers the failure the data actually exhibit. The alternative models compress the scale rather than shifting it, which Proposition 1 does not protect against and a standardised or rank-based differential does. The cross-model replication bears this out. Levels move by up to 9.1 index points and the raw gap roughly halves, three-quarters to nine-tenths of that movement being the affine rescaling Corollary 2 predicts, whilst the sign, the compositional anatomy and the scale-free magnitude of about two percentile points hold, the non-affine residual sitting at 23--24\% of the Cauchy--Schwarz worst case. Second, the compositional question posed by a devolved fiscal authority now has an answer, and the anatomy carries more of the information than the headline number does. Scotland shows a professional-management underweight and a care-sector overweight that persist even against the ex-London benchmark, with roughly three-quarters of the gap running within industries rather than through Scotland's sectoral structure. Third, the propagation template is itself a contribution. It puts a calibrated interval on the aggregate index, carries that interval through a TFP functional, checks it at the micro level with a wage gradient on which the scoring noise turns out not to bind, and closes with a stylised net-position translation under stated pass-through assumptions.

For a forecasting body the practical guidance is short. Use continuous indices in place of classified shares, because a classifier that codes 91\% of employment as exposed censors precisely the variation a regional comparison needs. Trust differentials over levels, because the differencing removes the model bias that makes a level a property of the scoring model. And expect the choice of model to move magnitudes, mostly by rescaling the index and only slightly by reordering it, which is why the scale-free forms of Corollary 2 should carry the reported result.

The limitations deserve stating plainly. Nothing here is causally identified, since no instrument for AI exposure yet matches the commuting-zone design of the automation literature, so what I report are compositional measurements. They do not identify displacement effects. Assumption 1, region invariance of within-occupation task content, cannot be tested without UK regional task data that no source provides, though the industry-margin diagnostics of \cref{app:data} bound the strain it is under. The productivity projection imports the Acemoglu calibration wholesale, so its levels inherit whatever that calibration gets wrong, which is why I rest on the differential in the implied paths. And the fiscal translation is an accounting exercise, not a forecast of the net position: its revenue elasticity is computed from the Scottish schedule rather than assumed, but the pass-through from productivity to earnings is imposed, and the population channel that dominates per-capita indexation is left out entirely.

What the framework offers extends beyond this application. Any exposure-based regional estimand a fiscal body builds, for any pair of regions and any scoring model, inherits the same common-mode structure: an employment-weighted contrast whose weights sum to one differences out whatever the model does uniformly, and the residual is bounded by the alignment of the model's bias with the share gap. The propagation template is reusable as the models change, which they will, and the calibration can be re-run against each new model at a cost far below that of the original scoring. The natural extensions are empirical. The ASHE earnings comparison under the sponsor's Data Access Agreement, the four-digit differential on the secure-access APS, and a longer panel for the first-difference specification the current window cannot support are all within reach.